% =============================================================================
% Front matter of the FeRRy paper: preamble, title, abstract, introduction,
% Table 1 and related works. Extracted on 8 Oct 2026 from Hermes_paper (6).tex
% (the live text only; commented and \iffalse'd drafts dropped), with a revised
% introduction. DeveloperDocs/paper/assemble.py appends methods.tex,
% evaluation.tex and back.tex to this file to build ferry_paper.tex.
% =============================================================================
\documentclass[conference]{IEEEtran}
\IEEEoverridecommandlockouts

% Text, Language, and URLs
\usepackage[utf8]{inputenc}
\usepackage[english]{babel}
\usepackage[hyphens]{url}
\usepackage{cite}

% Math and Symbols
\usepackage{amsmath,amssymb,amsfonts} % amssymb is included here
\usepackage{textcomp}

% Figures, Subfigures, and Colors
\usepackage{float}
\usepackage{graphicx}
\usepackage{subcaption}
\usepackage[table]{xcolor} % Updated to include [table] option

% Tables and Formatting
\usepackage{multirow}
\usepackage{tabularx}
\usepackage{booktabs} % Already existed
\usepackage{makecell} % Already existed
\usepackage{diagbox}
\usepackage{multicol}
\usepackage{fancyhdr}
\usepackage{colortbl} % Added for cell coloring
\usepackage{array}    % Added for advanced array/table options

% Custom Table Columns
\newcolumntype{L}{>{\raggedright\arraybackslash}X}
\newcolumntype{P}[1]{>{\centering\arraybackslash}m{#1}}

% Algorithms (Using algorithm + algpseudocode exclusively to avoid conflicts)
\usepackage{algorithm}
\usepackage{algpseudocode}

% TikZ and PGFPlots
\usepackage{tikz}
\usetikzlibrary{positioning, arrows.meta, calc}
\usepackage{pgfplots}
\pgfplotsset{compat=1.18} 

% Custom command definitions (defined exactly once)
\newcommand{\threat}[1]{\text{#1}}
\newcommand{\securitymechanism}[1]{\mathcal{S}_{\text{#1}}}
\newcommand{\protocol}[1]{\mathcal{C}_{\text{#1}}}
\newcommand{\experiment}[1]{\mathcal{E}_{\text{#1}}}
\newcommand{\performance}[1]{P_{\text{#1}}}

% FeRRy macros (methods and evaluation)
\newcommand{\sysname}{FeRRy}
\newcommand{\bbar}{\bar{b}}
\newcommand{\Tnom}{T_{\mathrm{nom}}}

\pagestyle{plain}

% \def\BibTeX{{\rm B\kern-.05em{\sc i\kern-.025em b}\kern-.08em
%     T\kern-.1667em\lower.7ex\hbox{E}\kern-.125emX}}


\makeatletter
\newcommand{\linebreakand}{%
  \end{@IEEEauthorhalign}
  \hfill\mbox{}\par
  \mbox{}\hfill\begin{@IEEEauthorhalign}
}
\makeatother

\begin{document}
% TODO(title): the paper now presents FeRRy; the old title named HERMES.
\title{FeRRy: RF-Aware Mission Planning for Federated Learning with Drones as Data Mules}
% Authors: omitted for double-anonymous review. The \author block (names,
% affiliations, emails) is kept in the Overleaf copy, not in this public repo;
% restore it there for the camera-ready.

\maketitle

% Abstract: the short hybrid of abstract_revision.md, adopted 8 Oct 2026
% (the document keeps the earlier versions and the audit of the original).
\begin{abstract}
Federated learning (FL) supports privacy by letting devices train a shared model
without sharing their raw data, but assumes every device can reach a central
server. In contested and disaster-zone deployments that assumption fails:
devices that are out of range or poorly connected are skipped repeatedly, and
federation rounds stall. Unmanned aerial vehicles (UAVs) acting as data mules can
carry model updates to and from such devices, but a mule has a finite mission
budget, and its radio decides how far it reaches. We present FeRRy, a
Federated RF-aware Routing framework that treats a mule's reach as a decision,
through three components. First, before each flight, a band-aware mission planner
chooses the contact band class, which trades range for rate, together with the
route, under one feasibility test of per-device deadlines and the mission budget,
with a coverage term that keeps devices from being starved. Second, in flight,
the mule re-plans the rest of its route when the observed channel uses up the
plan's slack, and a per-stop rule can switch to a faster band and reorder the
remaining stops, each reorder checked by the same test. Third, a hierarchical
aggregation layer merges updates on each mule, then across mules at the edge
server, weighting stale updates down. In a prototype whose server, mules and
devices run as separate processes over TCP and train a real intrusion-detection
model on CICIoT2023, with flight time, the radio channel and energy simulated, we
compare FeRRy with seven baseline schedulers, including adaptations of
mobile-transporter federated learning (FedEx), Oort, FedCS and a Whittle-index
age-of-update scheduler. With six devices, one mule and a moderate budget, FeRRy
reaches the target accuracy 27--31\% sooner than each baseline, because its plan
chooses a long-reach narrow band that removes most of the transit between stops.
Under tight budgets, it keeps the age of updates about a quarter lower than
age-of-information schedulers and, with 24 devices, reaches the target 12\%
sooner than FedEx's route. Its plans stay within a tight budget in 70\% of
missions, overrunning by only 8.5\,s on average otherwise.
\end{abstract}

\begin{IEEEkeywords}
Decentralized learning, federated learning, UAV data mules, tactical edge computing, mobile relay networks
\end{IEEEkeywords}

\section{Introduction}
\label{sec:introduction}

Disaster response, wildfire monitoring, agricultural sensing and tactical
operations increasingly rely on edge devices deployed where fixed infrastructure
is sparse, damaged or absent~\cite{zeng2016wireless, song2022comprehensive}.
These are also the settings most prone to denied, disrupted, intermittent and
limited (DDIL) connectivity~\cite{gupta2015survey}. Federated learning (FL) suits
them, because devices exchange model updates rather than raw
data~\cite{mcmahan2017communication, kostage2025hifins}, but it still needs every
device to reach the server: when a device cannot, its update is lost, and rounds
close short, time out, or fail to produce a usable model~\cite{li2020federated}.
The gap is not faster FL communication but federation that continues when the
network itself is fragmented.

UAVs acting as data mules can bridge disconnected regions by carrying model state
between devices and servers~\cite{zhang2025latency, arsalaan2024uavs}, but a mule
has limited flight time and energy, and a contact is useful only if an update is
ready and the link can sustain the transfer when the mule arrives. Hierarchical
FL, UAV-assisted FL, physical model transport, RF-adaptive networking and client
selection each address part of this problem~\cite{liu2020client, lin2023delay,
parasnis2023connectivity, yemini2023robust, mestoukirdi2022uav, bian2025indirect,
wang2018deep, cui2019multi, nishio2019client, lai2021oort}, but they retain
persistent reachability, stationary aggregation, or a communication setting that
the mobility decision does not control.

We present \sysname{}, a Federated RF-aware Routing framework whose central
observation is that a mule's radio reach is itself a decision. A contact band
with a narrower occupied bandwidth reaches farther at a lower rate: it can serve
a whole field from one stop at the cost of a longer dwell, whereas a wide band
serves each device quickly but only from nearby, so the mule spends most of its
mission in transit. Which trade wins depends on how many devices need service,
where they are and how much time remains, so \sysname{} chooses the band class
and the route together, before each flight.

\sysname{} decides on two clocks. On the \emph{plan clock}, at the dock, it
searches band classes and routes under one feasibility test of per-device
deadlines and the mission budget, and commits the plan that serves the most
coverage weight, a weight that grows with the missions since a device last
contributed. On the \emph{flight clock}, a per-stop rule can use the observed
channel to pick a faster band or reorder the remaining stops, and the mule
re-plans when the channel has used up the plan's slack. Updates are merged on
the mule with staleness weights and across mules at the server, so rounds close
with what arrives. We evaluate three variants: F commits the plan, FX adds a
fixed per-stop rule, and FQ a learned one.

In a multi-process prototype that trains a real intrusion-detection model on
CICIoT2023~\cite{CiCIoT2023}, with mission time, the channel and energy
simulated, we compare \sysname{} with seven whole schedulers under a
pre-registered analysis plan. With six devices and one mule, it reaches the
target accuracy 27--31\% sooner than each, because its plan chooses the
long-reach narrow band and cuts the transit between stops from about 140\,s to
15\,s; pinned to the wide band, it is no faster than FedEx's route. Under tight
budgets it keeps updates fresher than age-of-information schedulers. Learning
adds nothing to the fixed per-stop rule, and an exploratory comparison locates
the gain in the plan: inside it, the learned score reaches the target 36\%
sooner than a learned scheduler without one. We also report \sysname{}'s limits:
its plan prices the mean channel, so it overruns tight budgets in 30\% of
missions, and it does not price training time.

The contributions of this paper are:

\begin{itemize}

\item \textbf{Reach as a mission decision.} We treat the contact band class,
which sets how far a mule reaches and how fast it transfers, as a decision made
jointly with the route under a mission budget, and give a planner that searches
both under one feasibility test.

\item \textbf{A two-clock design with bounded learning.} The same test gates a
dock-time plan and in-flight per-stop decisions, with adaptive per-device
deadlines, a coverage weight and age cap against starvation, and
staleness-weighted merges; learning ranks only what the gates admit.

\item \textbf{A prototype and a simulator built from the same code.} A
multi-process prototype that trains a real model, and FerrySim, which trains the
learned components and plans for up to 96 devices.

\item \textbf{A pre-registered comparison with the state of the art.} Against
seven whole schedulers, on paired seeds with multiplicity correction, we report
where \sysname{} wins, why, which parts of its design have no measurable effect,
and where it falls short.

\end{itemize}

Section~\ref{sec:related_works} reviews related work,
Section~\ref{sec:method} presents \sysname{}'s design,
Section~\ref{sec:setup} the experimental setup and design,
Section~\ref{sec:results} the results, and Section~\ref{sec:conclusion}
concludes.

\definecolor{goodgreen}{RGB}{178,210,150}
\definecolor{ferrygreen}{RGB}{105,160,90}
\definecolor{badred}{RGB}{230,170,160}
\definecolor{nayellow}{RGB}{240,210,120}

\newcommand{\cmarkc}{\cellcolor{goodgreen}\textcolor{black!55}{\scriptsize$\checkmark$}}
\newcommand{\xmarkc}{\cellcolor{badred}\textcolor{black!40}{\scriptsize$\times$}}
\newcommand{\tmarkc}{\cellcolor{nayellow}\textcolor{black!65}{\scriptsize$\triangle$}}
\newcommand{\cmarkF}{\cellcolor{ferrygreen}\textcolor{white}{\textbf{\scriptsize$\checkmark$}}}

\begin{table*}[!t]
\caption{\small \textsc{Comparison of Related FL, UAV Mobility, and RF-Aware Coordination Approaches}}
\label{tab:ferry-comparison}
\centering
\resizebox{\textwidth}{!}{%
\begin{tabular}{l c c c c c c}
\toprule
\multicolumn{1}{c}{Framework}
&
\begin{tabular}[c]{@{}c@{}}Physical Model\\ Transport\end{tabular}
&
\begin{tabular}[c]{@{}c@{}}Explicit RF\\ Decision\end{tabular}
&
\begin{tabular}[c]{@{}c@{}}Mobility / Route\\ Decision\end{tabular}
&
\begin{tabular}[c]{@{}c@{}}Deadline /\\ Freshness Aware\end{tabular}
&
\begin{tabular}[c]{@{}c@{}}In-Mission\\ Replanning\end{tabular}
&
\begin{tabular}[c]{@{}c@{}}Joint FL--Mobility/RF\\ Coordination\end{tabular}
\\
\midrule

Hierarchical FL
{\scriptsize \cite{liu2020client}}
& \xmarkc & \xmarkc & \xmarkc & \xmarkc & \xmarkc & \xmarkc \\

Delay-Aware HFL
{\scriptsize \cite{lin2023delay}}
& \xmarkc & \xmarkc & \xmarkc & \tmarkc & \xmarkc & \xmarkc \\

FedCS
{\scriptsize \cite{nishio2019client}}
& \xmarkc & \xmarkc & \xmarkc & \cmarkc & \xmarkc & \xmarkc \\

Oort
{\scriptsize \cite{lai2021oort}}
& \xmarkc & \xmarkc & \xmarkc & \cmarkc & \xmarkc & \xmarkc \\

UAV Multi-Community FL
{\scriptsize \cite{mestoukirdi2022uav}}
& \xmarkc & \xmarkc & \cmarkc & \tmarkc & \xmarkc & \cmarkc \\

FedEx
{\scriptsize \cite{bian2025indirect}}
& \cmarkc & \xmarkc & \cmarkc & \tmarkc & \xmarkc & \cmarkc \\

Model-Aided Multi-UAV RL
{\scriptsize \cite{chen2023model}}
& \xmarkc & \tmarkc & \cmarkc & \xmarkc & \cmarkc & \xmarkc \\

Deadline-Sensitive Multi-UAV
{\scriptsize \cite{mondal2025deep}}
& \xmarkc & \tmarkc & \cmarkc & \cmarkc & \cmarkc & \xmarkc \\

DRL Multichannel Access
{\scriptsize \cite{wang2018deep}}
& \xmarkc & \cmarkc & \xmarkc & \xmarkc & \xmarkc & \xmarkc \\

MARL UAV Resource Allocation
{\scriptsize \cite{cui2019multi}}
& \xmarkc & \cmarkc & \tmarkc & \xmarkc & \xmarkc & \xmarkc \\

Value-Aware UAV FL
{\scriptsize \cite{cui2023data}}
& \xmarkc & \xmarkc & \xmarkc & \cmarkc & \xmarkc & \xmarkc \\

UAV-Enabled Asynchronous FL
{\scriptsize \cite{zhai2024uav}}
& \xmarkc & \tmarkc & \cmarkc & \cmarkc & \xmarkc & \cmarkc \\

FedBuff
{\scriptsize \cite{nguyen2022federated}}
& \xmarkc & \xmarkc & \xmarkc & \cmarkc & \xmarkc & \xmarkc \\

Async-HFL
{\scriptsize \cite{yu2023async}}
& \xmarkc & \xmarkc & \xmarkc & \cmarkc & \xmarkc & \xmarkc \\

\midrule

\textbf{FeRRy (ours)}
& \cmarkF
& \cmarkF
& \cmarkF
& \cmarkF
& \cmarkF
& \cmarkF \\

\bottomrule
\end{tabular}%
}

\vspace{2pt}
{\footnotesize
\cmarkc\ Present
\quad
\xmarkc\ Absent
\quad
\tmarkc\ Partial / indirect
}
\end{table*}


% Architecture figure: OUT of the build (8 Oct 2026). The drawing,
% HERMES_Case_Scenarios_compact.drawio, is the HERMES design: "WITH HERMES"
% labels, a DDQN choosing among carrier frequencies (FeRRy's band classes share
% one carrier and differ in bandwidth; nothing in FeRRy's RF layer is learned),
% and the old gated job scheduler. Redraw it to match Section III (plan clock,
% flight clock, merges), then uncomment this block and restore
% "A cluster (Fig.~\ref{fig:architecture})" in methods.tex's Setting paragraph.
% \begin{figure*}[t]
%     \centering
%     \includegraphics[trim=0cm 0cm 0cm 0cm, clip, scale=0.35]{Figures/HERMES_Case_Scenarios_compact.drawio.png}
%     \caption{\sysname{} overview: field devices, UAV data mules, and the edge server.}
%     \label{fig:architecture}
% \end{figure*}

\section{Related Works}
\label{sec:related_works}

Federated learning reduces centralized data movement by allowing participants to train locally and exchange model updates rather than raw datasets. FedAvg established this model-averaging paradigm, while FedProx improved robustness to statistical and systems heterogeneity \cite{mcmahan2017communication,li2020federated}. Hierarchical approaches subsequently introduced intermediate aggregation to reduce communication cost and accommodate heterogeneous edge environments. Aggregation-oriented hierarchical FL performs partial aggregation at edge servers before cloud synchronization, while delay-aware hierarchical FL explicitly accounts for communication delay between edge and cloud tiers \cite{liu2020client,lin2023delay}. Connectivity-aware semi-decentralized FL further exploits time-varying device-to-device links, and collaborative relaying allows well-connected clients to forward updates for clients experiencing intermittent server connectivity \cite{parasnis2023connectivity,yemini2023robust}. These approaches improve FL under unreliable communication, but they still rely on a communication graph in which model state is eventually forwarded through wireless links. FeRRy targets the more disruptive case in which such an end-to-end path may not exist and model state must instead be physically transported between isolated participants.

UAV-assisted FL introduces mobility into the learning infrastructure, but the role assigned to the UAV varies considerably. A UAV may operate as a mobile FL server whose routing is optimized with bandwidth, computation, and transmit power \cite{zhang2025latency}, or as a mobile orchestrator that visits discrete locations while optimizing routing and device scheduling altogether across multiple FL communities \cite{mestoukirdi2022uav}.  These systems use mobility to restore or improve communication between ground participants and the learning infrastructure. The closest architectural precedent to FeRRy is FedEx, which removes direct client-server connectivity and uses mobile transporters to physically carry model information \cite{bian2025indirect}. FedEx establishes physical model transport and relates FL performance to transporter travel time, but assumes preplanned tours and fixed communication conditions without adapting service to changing RF conditions, per-device deadlines, or in-mission contact opportunities. FeRRy extends this setting by treating reachability, service time, participant urgency, and communication conditions as part of the mission decision itself.

Another research direction studies UAV trajectory planning together with communication-aware data collection. Earlier formulations have demonstrated that trajectory and communication should not be optimized independently, and should jointly optimize multi-UAV trajectories, user scheduling, association, and transmit power for wireless service \cite{wu2018joint}, while more recent learning-based approaches move these decisions online. Multi-UAV controllers have used SNR, reachability, remaining device data, relative distance, and battery state to learn trajectories \cite{chen2023model}, and subsequent work has learned 3-D trajectories and device scheduling \cite{gao2024federated}, incorporated collection deadlines, link quality, and energy constraints through MAPPO \cite{mondal2025deep}, further optimized trajectory and task offloading under AoI, latency, and energy objectives \cite{wang2025joint}, and considered battery and interference constraints in federated-RL data collection \cite{kumar2026battery}. RF adaptation has likewise been studied through deep reinforcement learning for dynamic multichannel access \cite{wang2018deep} and multi-agent resource allocation over users, power levels, and subchannels in UAV networks \cite{cui2019multi}. These studies establish that communication-aware, and energy-aware multi-UAV trajectory learning is already well studied. FeRRy therefore does not claim RF-aware trajectory planning itself as novel, instead, its distinction is that the communication choice directly changes the feasibility and duration of transporting FL updates, while participant urgency determines which contact is worth serving next.

Client-selection and freshness-aware scheduling provides another part of this coordination problem. Existing methods select participants according to round deadlines and resource availability \cite{nishio2019client}, favor clients with higher training loss \cite{cho2022towards}, or use statistical utility and participation history observed in previous rounds \cite{lai2021oort}. This distinction between pre-selection and historical information is particularly important for physical data muling because isolated participants cannot continuously report their current state before the UAV reaches them. Related UAV-FL work has further balanced update age and data value under unstable connectivity \cite{cui2023data}, incorporated freshness into fairness-oriented scheduling \cite{zhu2024fairness}, required devices to be served within a bounded number of slots \cite{zhai2024uav}, and increased the priority of previously failed or unscheduled devices \cite{mestoukirdi2022uav}. These studies establish age-aware ranking and mechanisms to reduce starvation. FeRRy instead focuses on how these concerns interact with physical transport cost, where servicing an old or valuable update also requires flight time, contact time, sufficient radio reach, and remaining mission budget. 

Finally, asynchronous and hierarchical FL research addresses model-version mismatch when participant updates arrive at different times. Prior work has buffered asynchronous client updates before applying server-side aggregation, reducing sensitivity to individual stragglers while remaining compatible with secure aggregation \cite{nguyen2022federated}. Other hierarchical approaches perform asynchronous aggregation across both gateway and cloud tiers while jointly considering device selection and gateway association in heterogeneous IoT environments \cite{yu2023async}. These methods assume that updates are delivered through the network, whereas physical data muling introduces travel time as an additional source of staleness. FeRRy therefore treats aggregation as part of the mission lifecycle, where the same age and deadline information is used to determine whether a participant to be served also influences whether its transported update remains useful when the mule returns. The key distinction is thus the coupling of physical admission, service order, RF-dependent contact feasibility, and update freshness, rather than a new aggregation mechanism in isolation.

Table~\ref{tab:ferry-comparison} shows that prior work addresses individual parts of this problem, including FL-aware participant selection, freshness-aware scheduling, UAV trajectory planning, RF adaptation, and physical model transport. However, even the closest transporter-based approach does not address RF-dependent reach and service time with participant deadlines and in-mission route adaptation. FeRRy addresses this gap by using communication conditions, mobility, and FL-update urgency to plan and revise model transport under a finite mission budget.

